\usepackage{amsfonts,amssymb,amsmath,amsthm,mathtools,graphicx,bm,xspace,booktabs,array,multirow,dsfont}
\usepackage{enumitem}
\setlist[itemize]{noitemsep}
\setlist[enumerate]{noitemsep}
\usepackage[usenames,dvipsnames]{xcolor}
\definecolor{Gred}{RGB}{219,50,54}
\definecolor{Ggreen}{RGB}{60,186,84}
\definecolor{Gblue}{RGB}{72,133,237}
\definecolor{Gyellow}{RGB}{247,178,16}
\definecolor{ToCgreen}{RGB}{0,128,0}
\definecolor{myGold}{RGB}{231,141,20}
\definecolor{myBlue}{rgb}{0.19,0.41,.65}
\definecolor{myPurple}{RGB}{175,0,124}
\usepackage[colorlinks=true,hypertexnames=false]{hyperref}
\usepackage{tikz}
\usetikzlibrary{arrows.meta,shapes}
\usepackage[framemethod=tikz]{mdframed}
\mdfsetup{roundcorner=4pt,linewidth=.5,innerleftmargin=10pt,innerrightmargin=10pt,innertopmargin=10pt,innerbottommargin=10pt}
\usepackage{thmtools,thm-restate}
\usepackage[no-math]{fontspec}
\newtheorem{theorem}{Theorem}
\newtheorem{lemma}[theorem]{Lemma}
\newtheorem{corollary}[theorem]{Corollary}
\newtheorem{proposition}[theorem]{Proposition}
\newtheorem{claim}[theorem]{Claim}
\newtheorem{fact}[theorem]{Fact}
\newtheorem{techlemma}[theorem]{Technical Lemma}
\theoremstyle{definition}
\newtheorem{definition}[theorem]{Definition}
\theoremstyle{remark}
\newtheorem{remark}[theorem]{Remark}
\newtheorem{example}[theorem]{Example}
\newenvironment{Enumerate}{\begin{enumerate}}{\end{enumerate}}
\newenvironment{Itemize}{\begin{itemize}}{\end{itemize}}
\newcommand{\contour}[2]{#2}
\let\boldm\bm
\newenvironment{nproblem}[1][\unskip]{\medskip\begin{mdframed}\noindent\underline{$#1$ Problem.}\\\noindent}{\end{mdframed}}
\allowdisplaybreaks
\AtBeginDocument{\author{\parbox{.95\textwidth}{\centering Mika G\"o\"os, Pritish Kamath,\\ Katerina Sotiraki, Manolis Zampetakis}}}

\newcommand{\inparen}[1]{\left ( #1 \right )}
\newcommand{\infork}[1]{\left \{ \begin{matrix} #1 \end{matrix} \right .}
\newcommand{\inabs}[1]{\begin{vmatrix} #1 \end{vmatrix}}
\renewcommand{\mod}[1]{~\inparen{\mathrm{mod}~#1}}
\newcommand{\bx}{{\boldm x}}
\def\ceil#1{\mathop{\left\lceil#1\right\rceil}}
\def\abs#1{\left|#1\right|}
\def\p#1{\left(#1\right)}
\def\a#1{\langle#1\rangle}
\def\set#1{\left\{#1\right\}}
\newcommand{\parasc}[1]{\noindent \contour{black}{#1}}
\newcommand{\paragr}[1]{\noindent \textbf{\boldmath #1}}
\newcommand{\F}{\ensuremath{\mathbb{F}}}
\newcommand{\bbF}{\ensuremath{\mathbb{F}}}
\newcommand{\bbZ}{\ensuremath{\mathbb{Z}}}
\newcommand{\calC}{\mathcal{C}}
\newcommand{\calG}{\mathcal{G}}
\newcommand{\calM}{\mathcal{M}}
\newcommand{\calR}{\mathcal{R}}
\newcommand{\calT}{\mathcal{T}}
\newcommand{\calV}{\mathcal{V}}
\renewcommand{\vec}[1]{\boldsymbol{#1}}
\newcommand{\vecf}{\ensuremath{\boldsymbol{f}}}
\newcommand{\vecg}{\ensuremath{\boldsymbol{g}}}
\newcommand{\vech}{\ensuremath{\boldsymbol{h}}}
\newcommand{\vecv}{\ensuremath{\boldsymbol{v}}}
\newcommand{\vecx}{\ensuremath{\boldsymbol{x}}}
\newcommand{\vecy}{\ensuremath{\boldsymbol{y}}}
\newcommand{\vecz}{\ensuremath{\boldsymbol{z}}}
\newcommand{\bit}{\ensuremath{\set{0,1}}}
\DeclareMathOperator*{\poly}{\ensuremath\mathrm{poly}}
\newcommand{\problem}[1]{\ensuremath{\textsc{#1}}\xspace}
\newcommand{\class}[1]{\ensuremath{\mathsf{#1}}\xspace}
\newcommand{\FNP}{\class{FNP}}
\newcommand{\TFNP}{\class{TFNP}}
\newcommand{\PPA}{\class{PPA}}
\newcommand{\PPAp}{\class{PPA}_{{p}}}
\newcommand{\sis}{\problem{SIS}}
\newcommand{\pnum}{p}
\newcommand{\qnum}{q}
\newcommand{\lonelyp}{\ensuremath\textsc{Lonely}_{\pnum}}
\newcommand{\bipartitep}{\ensuremath{\textsc{Bipartite}_{\pnum}}}
\newcommand{\twoMatchingsp}{\ensuremath{\textsc{TwoMatchings}_{\pnum}}}
\newcommand{\sucBipartitep}{\ensuremath{\textsc{SuccinctBipartite}_{\pnum}}}
\newcommand{\leafq}{\ensuremath\textsc{Leaf}_{\qnum}}
\newcommand{\lonelyq}{\ensuremath\textsc{Lonely}_{\qnum}}
\newcommand{\bipartiteq}{\ensuremath{\textsc{Bipartite}_\qnum}}
\newcommand{\leaf}{\ensuremath\textsc{Leaf}}
\newcommand{\lonely}{\ensuremath\textsc{Lonely}}
\newcommand{\bipartite}{\ensuremath{\textsc{Bipartite}}}
\newcommand{\twoMatchings}{\ensuremath{\textsc{TwoMatchings}}}
\newcommand{\sucBipartite}{\ensuremath{\textsc{SuccinctBipartite}}}
\newcommand{\reducible}{\preceq}
\newcommand{\Chevalley}{\ensuremath\textsc{Chevalley}}
\newcommand{\CWT}{CWT\xspace}
\newcommand{\lChevalleyp}{\textsc{GeneralChevalley}_{\pnum}}
\newcommand{\cChevalleyp}{\textsc{ChevalleyWithSymmetry}_{\pnum}}
\newcommand{\eChevalleyp}{\textsc{Chevalley}_{\pnum}}
\newcommand{\lChevalley}{\textsc{GeneralChevalley}}
\newcommand{\cChevalley}{\textsc{ChevalleyWithSymmetry}}
\newcommand{\eChevalley}{\textsc{Chevalley}}
\newcommand{\CW}{\mathsf{CW}}

\providecommand{\bigtimes}{\mathop{\times}\displaylimits}
\providecommand{\bracks}[1]{\left[#1\right]}

\newcommand{\M}{\mathsf M}
\renewcommand{\P}{\mathsf P}
\newcommand{\amp}{\textsf{\upshape\&}}
\newcommand{\bamp}{~\textsf{\upshape\&}~}
\newcommand{\secref}[1]{\hyperref[#1]{\S\ref{#1}}}
\newcommand{\eqnref}[1]{\hyperref[#1]{Eq.~\ref{#1}}}
